Each criterion is recomputed for every retained conditional-quantile draw. Table entries report the posterior mean and equal-tailed 95\% credible interval of the resulting draw-wise criterion. Fit RMSE measures recovery of the true conditional quantile over the 500-observation training sample; forecast MAE and check loss average the 1,000 rolling-origin and horizon evaluations used in the point comparison. The intervals condition on the simulated data, evaluation design, and criterion-specific fitted model. Variational Bayes intervals are approximate.

The exDQLM entries use exdqlm 1.1.1; the DQLM and both Q--DESN entries are unchanged. The largest absolute change in an exDQLM posterior forecast-criterion mean is 0.8464\%; these forecast changes are not material, and the substantive conclusions remain stable. Point estimates and posterior intervals use distinct estimators: point estimates summarize a single fixed path for VB and chain-level fixed paths for MCMC, whereas interval centers are posterior means of draw-wise criteria. The two summaries are therefore interpreted separately.

5 of the 108 displayed MCMC summaries are marked by daggers to indicate diagnostic cautions; the supplement provides details.
